% 文件名：MainBody.tex
% 文件描述：四川大学2020研究生硕/博士 LaTeX 模版
% 作者：Kevin T. Lee [hello@lidengju.com]
% 修改日期：2020年6月1日
 
% 设置文档属性
% 参数说明
% professional: 专业学位
% academic: 学术学位
% master: 硕士
% doctor: 博士
% approval: 送审版本，将不生成声明
% secret: 保密论文，将显示密级
% color: 红色川大logo
% kfont=<⟨none|adobe|fandol|founder|mac|macnew|macold|ubuntu|windows|windowsnew|windowsold|...⟩>，
% 不填写则默认fandol
\documentclass[professional,master]{./Template/scuthesis2020}
% 各部分仍另页开始，但不为奇数页起始补空白页。
\let\cleardoublepage\clearpage
\begin{document}
% 设置文档信息
\unitid{10610} % 单位代码
\STUnumber{2024223045026} % 学号或送审编号
\securityClassification{秘密} % 密级：公开/内部/秘密/机密/绝密。当不使用secret时，不显示
\securityYear{3} % 保密年限
\CoverTitle{基于物理参数化高斯泼溅的} %封面标题
\CoverSubTitle{实时体积云渲染} % 无副标题时保留空参数，模板需要此项
\title{基于物理参数化高斯泼溅的实时体积云渲染} % 论文全称
\ENGtitle{Volumetric cloud rendering based on physically parameterized gaussian splatting} %论文全称英文
\school{计算机学院} % 培养单位
\ENGschool{College of Computer Science} % 培养单位英文
\author{向建鑫} % 作者姓名
\ENGauthor{Jianxin Xiang} % 作者英文名
\supervisor{李辉\quad 教授} % 指导教师
\ENGsupervisor{Prof. Hui Li} % 指导教师英文
\degreeclass{工程硕士} % 学位类别
\ENGdegreeclass{Master of Engineering} % 学位类别英文
\major{计算机技术} % 学科专业或领域名称
\ENGmajor{Computer Technology} % 学科专业或领域名称英文
\hasmajor{1} % 若有领域则为1，否则改为0
\defensedate{二〇二七年六月} % 论文答辩时间
\grantdate{二〇二七年六月} % 学位授予时间
\accomplishdate{二〇二七年三月} % 论文完成时间
\statementdate{June, 2027} % 声明时间
\direction{XXX}
\ENGdirection{Direction Name}
\keywords{研究生学位论文；写作指南；参考模板}
\ENGkeywords{graduate thesis;writing guide;reference template}

% 自动制作封面
\maketitle
% 自动制作中英文声明
\makestatement
% 设置论文正文前的页码、页眉等
\frontmatter\pagenumbering{Roman}\pagestyle{fancy}
% 包含摘要
% 保留中英文摘要页面，但不将摘要标题加入目录。
\begingroup
\renewcommand{\addcontentsline}[3]{}
\include{Chapters/0_0_Abstract}
\endgroup
% 自动制作目录
\maketoc
% 自动制作图表目录
\makefigtablist
% 设置论文正文部分的页码、页眉等
\mainmatter\pagenumbering{arabic}\pagestyle{fancy}
% 包含第一章、第二章等等
\include{Chapters/1_Introduction}
\include{Chapters/2_Structure}
\include{Chapters/3_Writing}
\include{Chapters/4_Printing}
\include{Chapters/5_Conclusion}

% 设置论文正文后的式样
\backmatter
% 按国标自动制作参考文献
% 参考文献数据文件为本目录下的ReferenceBase.bib
\begin{reference}
	\bibliography{ReferenceBase}
\end{reference}
% 包含在读期间科研成果
\include{Chapters/Achievements}
% 包含致谢
\makethanks
\end{document}
